\documentclass[10pt,a4paper]{amsart}
\usepackage{amsmath,amssymb,amsthm,graphicx}
% Retained preamble for the Stacks Project chapter snapshot in this directory.
% The chapter's own preamble.tex is a copy of the shared Stacks Project preamble; the block
% below is sliced MECHANICALLY, line for line, from it. Nothing was added, removed or
% redefined; the wrapper class, the xr-hyper/hyperref wiring and the externaldocument list
% are omitted because the assembler supplies the document class.
%
% lines taken from preamble.tex: 3, 5, 9, 11, 12, 156, 157, 158, ... (51 lines total)
%
% package coverage restored from the chapter's own preamble (2026-10-13): xy, xy-2cell

\usepackage[all]{xy}
\xyoption{2cell}
\UseAllTwocells

\newenvironment{reference}{\comment}{\endcomment}
\newenvironment{slogan}{\comment}{\endcomment}
\newenvironment{history}{\comment}{\endcomment}
\theoremstyle{plain}
\newtheorem{theorem}[subsection]{Theorem}
\newtheorem{proposition}[subsection]{Proposition}
\newtheorem{lemma}[subsection]{Lemma}
\theoremstyle{definition}
\newtheorem{definition}[subsection]{Definition}
\newtheorem{example}[subsection]{Example}
\newtheorem{exercise}[subsection]{Exercise}
\newtheorem{situation}[subsection]{Situation}
\theoremstyle{remark}
\newtheorem{remark}[subsection]{Remark}
\newtheorem{remarks}[subsection]{Remarks}
\numberwithin{equation}{subsection}
\def\lim{\mathop{\mathrm{lim}}\nolimits}
\def\colim{\mathop{\mathrm{colim}}\nolimits}
\def\Spec{\mathop{\mathrm{Spec}}}
\def\Hom{\mathop{\mathrm{Hom}}\nolimits}
\def\Ext{\mathop{\mathrm{Ext}}\nolimits}
\def\SheafHom{\mathop{\mathcal{H}\!\mathit{om}}\nolimits}
\def\SheafExt{\mathop{\mathcal{E}\!\mathit{xt}}\nolimits}
\def\Sch{\mathit{Sch}}
\def\Mor{\mathop{\mathrm{Mor}}\nolimits}
\def\Ob{\mathop{\mathrm{Ob}}\nolimits}
\def\Sh{\mathop{\mathit{Sh}}\nolimits}
\def\NL{\mathop{N\!L}\nolimits}
\def\CH{\mathop{\mathrm{CH}}\nolimits}
\def\proetale{{pro\text{-}\acute{e}tale}}
\def\etale{{\acute{e}tale}}
\def\QCoh{\mathit{QCoh}}
\def\Ker{\mathop{\mathrm{Ker}}}
\def\Im{\mathop{\mathrm{Im}}}
\def\Coker{\mathop{\mathrm{Coker}}}
\def\Coim{\mathop{\mathrm{Coim}}}
\DeclareMathSymbol{\boxtimes}{\mathbin}{AMSa}{"02}
\def\QCohstack{\mathcal{QC}\!\mathit{oh}}
\def\Cohstack{\mathcal{C}\!\mathit{oh}}
\def\Spacesstack{\mathcal{S}\!\mathit{paces}}
\def\Quotfunctor{\mathrm{Quot}}
\def\Hilbfunctor{\mathrm{Hilb}}
\def\Curvesstack{\mathcal{C}\!\mathit{urves}}
\def\Polarizedstack{\mathcal{P}\!\mathit{olarized}}
\def\Complexesstack{\mathcal{C}\!\mathit{omplexes}}
\def\Pic{\mathop{\mathrm{Pic}}\nolimits}
\def\Picardstack{\mathcal{P}\!\mathit{ic}}
\def\Picardfunctor{\mathrm{Pic}}
\def\Deformationcategory{\mathcal{D}\!\mathit{ef}}
\renewenvironment{reference}{}{}
\renewenvironment{slogan}{}{}
\renewenvironment{history}{}{}
% empty out citations inside navigation-pointer environments (no bibliography
% backend; spans themselves are cite-free - scanned 2026-10-14, lane-s3)
\renewcommand{\cite}[2][]{}
\usepackage{cleveref}
\crefname{originallemma}{original lemma}{original lemmas}
\Crefname{originallemma}{Original lemma}{Original lemmas}
\makeatletter
\newlabel{algebra-lemma-characterize-dvr}{{00PD}{source}{}{sourcepointer-1}{}}
\newlabel{algebra-lemma-characterize-dvr@cref}{{[originallemma][0][]00PD}{[1][0][]source}{}{}{}}
\makeatother
\crefname{originallemma}{original lemma}{original lemmas}
\Crefname{originallemma}{Original lemma}{Original lemmas}
\makeatletter
\newlabel{algebra-lemma-dimension-base-fibre-equals-total}{{00ON}{source}{}{sourcepointer-2}{}}
\newlabel{algebra-lemma-dimension-base-fibre-equals-total@cref}{{[originallemma][0][]00ON}{[1][0][]source}{}{}{}}
\makeatother
\crefname{originallemma}{original lemma}{original lemmas}
\Crefname{originallemma}{Original lemma}{Original lemmas}
\makeatletter
\newlabel{algebra-lemma-localization-of-regular-local-is-regular}{{0AFS}{source}{}{sourcepointer-3}{}}
\newlabel{algebra-lemma-localization-of-regular-local-is-regular@cref}{{[originallemma][0][]0AFS}{[1][0][]source}{}{}{}}
\makeatother
\crefname{originallemma}{original lemma}{original lemmas}
\Crefname{originallemma}{Original lemma}{Original lemmas}
\makeatletter
\newlabel{algebra-lemma-regular-normal}{{0567}{source}{}{sourcepointer-4}{}}
\newlabel{algebra-lemma-regular-normal@cref}{{[originallemma][0][]0567}{[1][0][]source}{}{}{}}
\makeatother
\crefname{originallemma}{original lemma}{original lemmas}
\Crefname{originallemma}{Original lemma}{Original lemmas}
\makeatletter
\newlabel{algebra-proposition-going-down-normal-integral}{{00H8}{source}{}{sourcepointer-5}{}}
\newlabel{algebra-proposition-going-down-normal-integral@cref}{{[originallemma][0][]00H8}{[1][0][]source}{}{}{}}
\makeatother
\crefname{originallemma}{original lemma}{original lemmas}
\Crefname{originallemma}{Original lemma}{Original lemmas}
\makeatletter
\newlabel{descent-lemma-descending-property-etale}{{02VN}{source}{}{sourcepointer-6}{}}
\newlabel{descent-lemma-descending-property-etale@cref}{{[originallemma][0][]02VN}{[1][0][]source}{}{}{}}
\makeatother
\crefname{originallemma}{original lemma}{original lemmas}
\Crefname{originallemma}{Original lemma}{Original lemmas}
\makeatletter
\newlabel{descent-lemma-etale-etale-local-source}{{036W}{source}{}{sourcepointer-7}{}}
\newlabel{descent-lemma-etale-etale-local-source@cref}{{[originallemma][0][]036W}{[1][0][]source}{}{}{}}
\makeatother
\crefname{originallemma}{original lemma}{original lemmas}
\Crefname{originallemma}{Original lemma}{Original lemmas}
\makeatletter
\newlabel{etale-section-properties-permanence}{{025L}{source}{}{sourcepointer-8}{}}
\newlabel{etale-section-properties-permanence@cref}{{[originallemma][0][]025L}{[1][0][]source}{}{}{}}
\makeatother
\crefname{originallemma}{original lemma}{original lemmas}
\Crefname{originallemma}{Original lemma}{Original lemmas}
\makeatletter
\newlabel{lemma-local-purity}{{0BMA}{source}{}{sourcepointer-9}{}}
\newlabel{lemma-local-purity@cref}{{[originallemma][0][]0BMA}{[1][0][]source}{}{}{}}
\makeatother
\crefname{originallemma}{original lemma}{original lemmas}
\Crefname{originallemma}{Original lemma}{Original lemmas}
\makeatletter
\newlabel{more-algebra-section-permanence-etale}{{0AGY}{source}{}{sourcepointer-10}{}}
\newlabel{more-algebra-section-permanence-etale@cref}{{[originallemma][0][]0AGY}{[1][0][]source}{}{}{}}
\makeatother
\crefname{originallemma}{original lemma}{original lemmas}
\Crefname{originallemma}{Original lemma}{Original lemmas}
\makeatletter
\newlabel{more-morphisms-lemma-etale-makes-quasi-finite-finite-at-point}{{02LK}{source}{}{sourcepointer-11}{}}
\newlabel{more-morphisms-lemma-etale-makes-quasi-finite-finite-at-point@cref}{{[originallemma][0][]02LK}{[1][0][]source}{}{}{}}
\makeatother
\title{Purity of branch locus (Nagata): a quasi-finite morphism f : X -> Y of locally Noetherian schemes, with O\_\{X,x\} normal, O\_\{Y,y\} regular, equal relative dimensions at least 1 at x, and f unramified in codimension 1 around x, is etale at x.}
\author{The Stacks Project Authors}
\date{stacks-project:pione.tex:local-snapshot-20261012; GNU FDL 1.2 or later}
\begin{document}
\maketitle
\noindent\textit{Editorial note.} Statement span is lines 5334-5365 of sources/stacks-pione/pinned-original-source.tex; proof span is lines 5367-5430, the complete author proof (62 body lines).\par\medskip
\noindent\textit{Editorial note.} Navigation stubs: 11 cited-but-not-included labels are supplied as non-load-bearing pointers with their pinned tags where the snapshot tag index has them.\par\medskip
\noindent\textit{Editorial note.} extra\_packages is empty on purpose: the retained preamble already loads xy with options and a second declaration clashes in the assembler wrapper.\par\medskip
\noindent\textit{Editorial note.} Text is retained verbatim, including the author's own typos and placeholders.\par\medskip
\noindent\textit{Editorial note.} Queue index 16 of queue\_pione.json (pb=2947, stmt 5334-5365, proof 5367-5430).\par\medskip
\noindent\textit{Editorial note.} pione wave-1: fresh-adjudication accept (2026-10-14); label not covered by the 85 old-session S3-pi cards.\par\medskip
\noindent\textit{Allowed standard prerequisites (editorial).}\begin{itemize}
\item Explicit prerequisite boundary. The cited inputs are prerequisites; the argument in the retained span is not replaced by any cited result.
\end{itemize}\medskip
\ifdefined\hypertarget\hypertarget{sourcepointer-1}{}\fi
\noindent\textit{Source reference.} Original lemma 00PD: sources/stacks-pione/pinned-original-source.tex, source lines 5430--5430. Cited by the retained statement or proof but not part of the collected spans; supplied only as a navigation pointer. The argument does not depend on its content.\par\medskip
\ifdefined\hypertarget\hypertarget{sourcepointer-2}{}\fi
\noindent\textit{Source reference.} Original lemma 00ON: sources/stacks-pione/pinned-original-source.tex, source lines 5430--5430. Cited by the retained statement or proof but not part of the collected spans; supplied only as a navigation pointer. The argument does not depend on its content.\par\medskip
\ifdefined\hypertarget\hypertarget{sourcepointer-3}{}\fi
\noindent\textit{Source reference.} Original lemma 0AFS: sources/stacks-pione/pinned-original-source.tex, source lines 5430--5430. Cited by the retained statement or proof but not part of the collected spans; supplied only as a navigation pointer. The argument does not depend on its content.\par\medskip
\ifdefined\hypertarget\hypertarget{sourcepointer-4}{}\fi
\noindent\textit{Source reference.} Original lemma 0567: sources/stacks-pione/pinned-original-source.tex, source lines 5430--5430. Cited by the retained statement or proof but not part of the collected spans; supplied only as a navigation pointer. The argument does not depend on its content.\par\medskip
\ifdefined\hypertarget\hypertarget{sourcepointer-5}{}\fi
\noindent\textit{Source reference.} Original lemma 00H8: sources/stacks-pione/pinned-original-source.tex, source lines 5430--5430. Cited by the retained statement or proof but not part of the collected spans; supplied only as a navigation pointer. The argument does not depend on its content.\par\medskip
\ifdefined\hypertarget\hypertarget{sourcepointer-6}{}\fi
\noindent\textit{Source reference.} Original lemma 02VN: sources/stacks-pione/pinned-original-source.tex, source lines 5430--5430. Cited by the retained statement or proof but not part of the collected spans; supplied only as a navigation pointer. The argument does not depend on its content.\par\medskip
\ifdefined\hypertarget\hypertarget{sourcepointer-7}{}\fi
\noindent\textit{Source reference.} Original lemma 036W: sources/stacks-pione/pinned-original-source.tex, source lines 5430--5430. Cited by the retained statement or proof but not part of the collected spans; supplied only as a navigation pointer. The argument does not depend on its content.\par\medskip
\ifdefined\hypertarget\hypertarget{sourcepointer-8}{}\fi
\noindent\textit{Source reference.} Original lemma 025L: sources/stacks-pione/pinned-original-source.tex, source lines 5430--5430. Cited by the retained statement or proof but not part of the collected spans; supplied only as a navigation pointer. The argument does not depend on its content.\par\medskip
\ifdefined\hypertarget\hypertarget{sourcepointer-9}{}\fi
\noindent\textit{Source reference.} Original lemma 0BMA: sources/stacks-pione/pinned-original-source.tex, source lines 5284--5292. Cited by the retained statement or proof but not part of the collected spans; supplied only as a navigation pointer. The argument does not depend on its content.\par\medskip
\ifdefined\hypertarget\hypertarget{sourcepointer-10}{}\fi
\noindent\textit{Source reference.} Original lemma 0AGY: sources/stacks-pione/pinned-original-source.tex, source lines 5430--5430. Cited by the retained statement or proof but not part of the collected spans; supplied only as a navigation pointer. The argument does not depend on its content.\par\medskip
\ifdefined\hypertarget\hypertarget{sourcepointer-11}{}\fi
\noindent\textit{Source reference.} Original lemma 02LK: sources/stacks-pione/pinned-original-source.tex, source lines 5430--5430. Cited by the retained statement or proof but not part of the collected spans; supplied only as a navigation pointer. The argument does not depend on its content.\par\medskip
\begin{lemma}[Purity of branch locus]
\label{lemma-purity}
\begin{reference}
\cite{Nagata-Purity} and \cite[Exp. X, Thm. 3.1]{SGA1}
\end{reference}
\begin{history}
This result was first stated and proved by Zariski in
geometric form in \cite{Zariski-Purity}.
The generalization to nonperfect ground fields by Nagata
was published as the next article in the same volume of the
Proceedings of the National Academy of Sciences of the United States of America
in \cite{Nagata-Remarks-Purity}. In the following year Nagata
proved the result for Noetherian local rings in \cite{Nagata-Purity}.
His proof uses a result of Chow which is a Bertini theorem for
complete local domains, see \cite{Chow-Bertini};
the history of Bertini's theorems is discussed in
Kleiman's historical article \cite{Kleiman-Bertini}.
A few years later a completely different proof was found by
Auslander, see \cite{Auslander-Purity}.
\end{history}
Let $f : X \to Y$ be a morphism of locally Noetherian schemes.
Let $x \in X$ and set $y = f(x)$. Assume
\begin{enumerate}
\item $\mathcal{O}_{X, x}$ is normal,
\item $\mathcal{O}_{Y, y}$ is regular,
\item $f$ is quasi-finite at $x$,
\item $\dim(\mathcal{O}_{X, x}) = \dim(\mathcal{O}_{Y, y}) \geq 1$
\item for specializations $x' \leadsto x$ with
$\dim(\mathcal{O}_{X, x'}) = 1$ our $f$ is unramified at $x'$.
\end{enumerate}
Then $f$ is \'etale at $x$.
\end{lemma}

\begin{proof}
We will prove the lemma by induction on
$d = \dim(\mathcal{O}_{X, x}) = \dim(\mathcal{O}_{Y, y})$.

\medskip\noindent
An uninteresting case is when $d = 1$.
In that case we are assuming that $f$ is unramified at $x$
and that $\mathcal{O}_{Y, y}$ is a discrete valuation ring
(Algebra, Lemma \ref{algebra-lemma-characterize-dvr}).
Then $\mathcal{O}_{X, x}$ is flat over $\mathcal{O}_{Y, y}$
(otherwise the map would not be quasi-finite at $x$)
and we see that $f$ is flat at $x$. Since flat $+$
unramified is \'etale we conclude (some details omitted).

\medskip\noindent
The case $d \geq 2$. We will use induction on $d$ to reduce
to the case discussed in Lemma \ref{lemma-local-purity}.
To check $f$ is \'etale at $x$ we may work \'etale locally
on the base and on the target
(Descent, Lemmas \ref{descent-lemma-descending-property-etale} and
\ref{descent-lemma-etale-etale-local-source}).
Thus we can apply More on Morphisms, Lemma
\ref{more-morphisms-lemma-etale-makes-quasi-finite-finite-at-point}
and assume that $f : X \to Y$ is finite and that $x$ is the unique
point of $X$ lying over $y$. Here we use that \'etale extensions of
local rings do not change dimension, normality, and regularity, see
More on Algebra, Section \ref{more-algebra-section-permanence-etale}
and
\'Etale Morphisms, Section \ref{etale-section-properties-permanence}.

\medskip\noindent
Next, we can base change by $\Spec(\mathcal{O}_{Y, y})$
and assume that $Y$ is the spectrum of a regular local ring.
It follows that $X = \Spec(\mathcal{O}_{X, x})$ as
every point of $X$ necessarily specializes to $x$.

\medskip\noindent
The ring map $\mathcal{O}_{Y, y} \to \mathcal{O}_{X, x}$ is
finite and necessarily injective (by equality of dimensions).
We conclude we have going down for
$\mathcal{O}_{Y, y} \to \mathcal{O}_{X, x}$ by
Algebra, Proposition \ref{algebra-proposition-going-down-normal-integral}
(and the fact that a regular ring is a normal ring by
Algebra, Lemma \ref{algebra-lemma-regular-normal}).
Pick $x' \in X$, $x' \not = x$ with image $y' = f(x')$.
Then $\mathcal{O}_{X, x'}$ is normal as a localization
of a normal domain. Similarly, $\mathcal{O}_{Y, y'}$ is
regular (see Algebra, Lemma
\ref{algebra-lemma-localization-of-regular-local-is-regular}).
We have $\dim(\mathcal{O}_{X, x'}) = \dim(\mathcal{O}_{Y, y'})$ by
Algebra, Lemma \ref{algebra-lemma-dimension-base-fibre-equals-total}
(we checked going down above).
Of course these dimensions are strictly less than $d$ as $x' \not = x$
and by induction on $d$ we conclude that $f$ is \'etale at $x'$.

\medskip\noindent
Thus we arrive at the following situation: We have a finite
local homomorphism $A \to B$ of Noetherian local rings
of dimension $d \geq 2$, with $A$ regular, $B$ normal, which
induces a finite \'etale morphism $V \to U$ on punctured spectra.
Our goal is to show that $A \to B$ is \'etale.
This follows from Lemma \ref{lemma-local-purity}
and the proof is complete.
\end{proof}

\end{document}
